% ch09_conclusion.tex : Conclusion, Limitations, and Future Work
% Dr. Yael Mizrahi : LaTeX Technical Author
% XeLaTeX + bidi + fontspec

\selectlanguage{hebrew}

\section{סיכום, מגבלות ועבודה עתידית}
\label{sec:ch09_conclusion}

\subsection{סיכום התרומות}
\label{sec:ch09_summary}

מאמר זה הציג מערכת ניווט ביומימטית לרחפנים זעירים בהשראת אקולוקציית עטלפים ואינטגרציה רב-חושית. המערכת משלבת מדידות טווח-דופלר מסונאר אולטראסוני, נתוני \en{IMU} וזרימה אופטית, תוך שימוש במסגרת מיזוג חיישנים רב-מודאלית המחקה את האינטגרציה הרב-חושית בקוליקולוס העליון של העטלף.

התרומות העיקריות של עבודה זו הן שלוש. ראשית, מסגרת מיזוג חיישנים ביומימטית המשלבת מדידות טווח-דופלר מסונאר, נתוני \en{IMU} וזרימה אופטית, בהשראת האינטגרציה הרב-חושית בקוליקולוס העליון של העטלף. שנית, מסנן קלמן מורחב (\en{EKF}) מודע-דופלר המנצל מדידות מהירות רדיאלית מהדים סונאריים לשיפור אמידת תנועה עצמית. שלישית, אימות ניסויי על רחפן \en{Crazyflie 2.1} בסביבות פנימיות נטולות \en{GPS}, המדגים שיפור של 40 אחוזים בדיוק אמידת המהירות.

תוצאות הסימולציה והניסויים מאמתות את היתרונות של הגישה המוצעת. ה-\en{EKF} מודע-הדופלר משיג \en{RMSE} מיקום של 0.12 מטרים ו-\en{RMSE} מהירות של 0.08 מטרים לשנייה, שיפור של 33 אחוזים ו-43 אחוזים בהתאמה בהשוואה ל-\en{EKF} הסטנדרטי. \en{FastSLAM 2.0} עם סיוע דופלר מאפשר מיפוי עקבי עם שיעור הצלחת סגירת לולאה של 78 אחוזים. המערכת שומרת על לוקליזציה יציבה גם בנוכחות כשלי חיישנים, עם התדרדרות בחן.

ההקבלה בין ניווט תת-ימי לניווט רחפנים בתת-קרקע מספקת תוקף בין-תחומי לגישה המוצעת. מערכת \en{DVL-aided INS} הוכחה באלפי פריסות \en{AUV}, והשיגה שגיאות מיקום של 0.5 עד 2 מטרים במשימות של 10 דקות. אותה מסגרת אלגוריתמית, המותאמת לאקוסטיקה אווירית ולדינמיקת רחפנים זעירים, מהווה את ליבת מערכת הניווט המוצעת.

\begin{equation}
\text{CRLB}(r) = \frac{c^2}{8 \pi^2 B^2 \text{SNR}}
\label{eq:ch09_crlb}
\end{equation}

\begin{equation}
\mathbf{u}_k^* = \arg\max_{\mathbf{u}_k} I(\mathbf{x}_k; \mathbf{z}_k | \mathbf{u}_k, \mathbf{z}_{1:k-1})
\label{eq:ch09_sensor_scheduling}
\end{equation}

\subsection{מגבלות}
\label{sec:ch09_limitations}

למרות התוצאות המעודדות, למערכת המוצעת מספר מגבלות משמעותיות.

\textbf{טווח סונאר מוגבל.} טווח הסונאר מוגבל ל-5 מטרים בשל הנחתה אקוסטית באוויר והספק השיא המוגבל של המתמרים. מגבלה זו מצמצמת את תחולת המערכת לסביבות פנימיות או קרובות-מכשולים. בסביבות חיצוניות פתוחות, טווח זה אינו מספק לניווט בטוח.

\textbf{ביצועים בחללים מהדהדים.} ביצועי המערכת מתדרדרים בחללים מהדהדים מאוד (למשל, חדרי בטון עם משטחים קשים). הדים מרובים ממשטחים שונים יוצרים עמימות בזיהוי ההד הראשוני, ועלולים לגרום לשגיאות אמידת טווח. מערכת ביטול ההדים הנוכחית אינה מספקת דיכוי מלא של הדים מרובים.

\textbf{קנה מידה של מסנן החלקיקים.} גישת מסנן החלקיקים (\en{FastSLAM 2.0}) מתרחבת בצורה גרועה עם גודל המפה. מספר החלקיקים הנדרש לשמירה על דיוק אמידה גדל ליניארית עם מספר נקודות הציון, מה שמגביל את התחולה לסביבות קטנות מ-1,000 מטרים רבועים.

\textbf{תלות בצורת גל ידועה.} המערכת מניחה שצורת הגל המשודרת ידועה בדיוק. שינויים בתגובת התדר של המתמר ובמהירות הקול התלויה בטמפרטורה עלולים לפגוע בביצועי ה-\en{matched filtering}. כיול טמפרטורה אוטומטי נדרש לפעולה אמינה בתנאי סביבה משתנים.

\textbf{צימוד אקוסטי בביטול רעש.} מסנן ה-\en{NLMS} דורש אות ייחוס המתואם עם רעש המדחף אך לא עם ההד. בפועל, מיקרופון משני המוצב ליד המנוע מספק ייחוס זה, אך צימוד אקוסטי בין המיקרופון הראשי למשני עלול לגרום לביטול אות חלקי של ההד.

\begin{table}[htbp]
\centering
\caption{סיכום מגבלות המערכת ופתרונות עתידיים מוצעים. (Summary of system limitations and proposed future solutions.)}
\label{tab:ch09_limitations}
\adjustbox{max width=\columnwidth}{%
\begin{tabular}{lp{6cm}p{6cm}}
\toprule
\textbf{מגבלה} & \textbf{תיאור} & \textbf{פתרון עתידי} \\
\midrule
טווח סונאר מוגבל & 5 מטרים מקסימום & מתמרים בעלי רגישות גבוהה יותר, אפנון אדפטיבי \\
הדהוד בחללים סגורים & ביצועים מתדרדרים & עיבוד הדים מרובים, מודל אקוסטי של החלל \\
קנה מידה של \en{FastSLAM} & מוגבל ל-1,000 מ"ר & \en{GraphSLAM} מבוזר, מיפוי היררכי \\
תלות בצורת גל & רגיש לטמפרטורה & כיול אוטומטי, צורות גל אדפטיביות \\
צימוד אקוסטי & ביטול אות חלקי & מיקרופון ייחוס מבודד, ביטול רעש ספקטרלי \\
\bottomrule
\end{tabular}%
}
\end{table}

\subsection{כיוונים עתידיים}
\label{sec:ch09_future_work}

\subsubsection{בחירה אדפטיבית של צורת גל סונאר}
\label{sec:ch09_adaptive_waveform}

עבודה עתידית כוללת בחירה אדפטיבית של צורת גל סונאר בהשראת מעבר \en{CF-FM} בעטלפים. עטלפים עוברים בין פולסי \en{CF} (אמידת דופלר מדויקת) לפולסי \en{FM} (רזולוציית טווח גבוהה) בהתאם לתנאי הסביבה. אסטרטגיה אדפטיבית זו ניתנת למימוש במערכת הסונאר כדי למטב בין דיוק דופלר לרזולוציית טווח בהתאם למשימת הניווט.

ניתן למסגר בעיה זו כבעיית תזמון אופטימלי תחת אי-ודאות, תוך שימוש בחסם \en{Cramér-Rao} (משוואה~\ref{eq:ch09_crlb}) כמדד לביצועים. מטרת התזמון האופטימלי היא לבחור את צורת הגל המספקת את תוספת המידע הגדולה ביותר ביחס לעלות ההספק שלה, כפי שמתואר במשוואה~\ref{eq:ch09_sensor_scheduling}.

\subsubsection{למידת חיזוק עמוקה לתזמון חיישנים}
\label{sec:ch09_rl_scheduling}

למידת חיזוק עמוקה לתזמון חיישנים תחת אי-ודאות היא כיוון מבטיח נוסף. סוכן למידת חיזוק יכול ללמוד מדיניות אופטימלית לבחירת אופן החישה בכל צעד זמן, תוך איזון בין דיוק אמידה, צריכת הספק ועומס חישובי. פרמטר ה-$\omega$ באלגוריתם \en{Covariance Intersection} יכול להילמד על ידי רשת עצבית המקבלת את הקווריאנסים של האמידות כקלט.

\subsubsection{שילוב עם זיהוי מקום חזותי}
\label{sec:ch09_visual_place}

שילוב עם זיהוי מקום חזותי (\en{visual place recognition}) ל-\en{SLAM} בקנה מידה גדול יותר יאפשר למערכת לנווט בסביבות נרחבות יותר. שימוש בתכונות \en{DINOv2 SALAD} לזיהוי מקום חזותי משיג \en{Recall@1} של 95.7 אחוזים על ערכת הנתונים \en{Pitts30k} \cite{CrewAIDocs}. שילוב של סגירת לולאה חזותית ואקוסטית עם אימות צולב יספק דיוק גבוה עוד יותר.

\subsubsection{SLAM שיתופי עם רחפנים מרובים}
\label{sec:ch09_multi_drone}

\en{SLAM} שיתופי עם רחפנים מרובים באמצעות תקשורת סונאר הוא כיוון מבטיח נוסף. רחפנים מרובים יכולים לחלוק מפות סונאר ולשפר את הדיוק והכיסוי המרחבי. תקשורת אקוסטית בין רחפנים, בהשראת תקשורת עטלפים, תאפשר תיאום ללא צורך בתשתית תקשורת נוספת.

\subsubsection{מימוש נוירומורפי מלא}
\label{sec:ch09_neuromorphic}

מימוש נוירומורפי מלא של כל צינור עיבוד האותות על שבב \en{Intel Loihi} יאפשר הפחתה דרמטית נוספת בצריכת ההספק. ה-\en{SNN} המתואר בפרק 7 מהווה צעד ראשון בכיוון זה. מימוש מלא של ה-\en{EKF} וה-\en{FastSLAM} בארכיטקטורה נוירומורפית ידרוש פיתוח של נוירונים המבצעים פעולות מטריציאליות, אך הפוטנציאל לחיסכון בהספק הוא משמעותי.

\subsection{סיכום}
\label{sec:ch09_final_summary}

מאמר זה הדגים את היתכנותה של גישה ביומימטית לניווט רחפנים זעירים בסביבות נטולות \en{GPS}. ההשראה מהאקולוקציה והאינטגרציה הרב-חושית של העטלף הובילה לפיתוח מערכת המשלבת סונאר, \en{IMU} וזרימה אופטית במסגרת מיזוג חיישנים אחידה. התוצאות מראות שיפור משמעותי בביצועי הניווט, במיוחד באמידת מהירות, הודות לשילוב מדידות דופלר מסונאר.

המערכת המוצעת מהווה צעד חשוב לקראת רחפנים אוטונומיים המסוגלים לנווט בסביבות מורכבות ללא תלות ב-\en{GPS}. השילוב של חיישנים בעלי הספק נמוך, אלגוריתמי מיזוג יעילים, וארכיטקטורת תוכנה מודולרית מאפשר פריסה על רחפנים זעירים במשקל של מאות גרמים בודדים. עבודה עתידית תתמקד בהרחבת טווח הפעולה, שיפור העמידות בתנאי סביבה קשים, ומימוש נוירומורפי מלא להפחתת צריכת ההספק נוספת \cite{Ulanovsky2008BatBrain, Yovel2010OptimalLocalization, Vanderelst2016PlaceRecognition}.

\subsection{ניתוח השפעה רחב}
\label{sec:ch09_broader_impact}

לגישה המוצעת יש השלכות רחבות מעבר לניווט רחפנים זעירים. היכולת לנווט בסביבות נטולות \en{GPS} רלוונטית למגוון יישומים: חיפוש והצלה במבנים שקרסו, מיפוי מערות ומנהרות, ניטור תשתיות תת-קרקעיות, וחקלאות מדויקת בחופות יער. השילוב של חיישנים בעלי הספק נמוך עם אלגוריתמי מיזוג יעילים מאפשר פריסה על פלטפורמות זעירות וזולות, מה שמנגיש טכנולוגיית ניווט אוטונומית למגוון רחב של משתמשים.

מנקודת מבט ביולוגית, העבודה מאששת השערות לגבי מנגנוני הניווט של עטלפים. היכולת לשחזר ביצועים דומים לאלה של עטלפים באמצעות מערכת מהונדסת פשוטה יחסית תומכת בהשערה שהעקרונות האלגוריתמיים הבסיסיים של אקולוקציה ניתנים למימוש באמצעות עיבוד אותות סטנדרטי. עם זאת, המורכבות העצומה של המוח העטלפי, הכוללת אזורים ייעודיים לעיבוד אקוסטי, מרחבי וזמני, עדיין רחוקה ממימוש מלא במערכות מהונדסות.

\appendix
\section{רשימת סמלים ומשתנים}
\label{sec:ch09_appendix}

\begin{table}[htbp]
\centering
\caption{רשימת סמלים ומשתנים בהם נעשה שימוש במאמר. (List of symbols and variables used throughout the paper.)}
\label{tab:ch09_symbols}
\adjustbox{max width=\columnwidth}{%
\begin{tabular}{lll}
\toprule
\textbf{סמל} & \textbf{הגדרה} & \textbf{יחידות} \\
\midrule
$\mathbf{x}_k$ & וקטור מצב בזמן $k$ & משתנה \\
$\mathbf{u}_k$ & וקטור בקרה בזמן $k$ & משתנה \\
$\mathbf{z}_k$ & וקטור מדידה בזמן $k$ & משתנה \\
$\mathbf{F}_k$ & מטריצת מערכת בזמן $k$ & חסרת ממדים \\
$\mathbf{H}_k$ & מטריצת מדידה בזמן $k$ & חסרת ממדים \\
$\mathbf{Q}_k$ & קווריאנס רעש תהליך & משתנה \\
$\mathbf{R}_k$ & קווריאנס רעש מדידה & משתנה \\
$\mathbf{P}_k$ & קווריאנס אמידת מצב & משתנה \\
$\mathbf{K}_k$ & רווח קלמן בזמן $k$ & חסר ממדים \\
$c$ & מהירות הקול באוויר & 343 מטרים/שנייה \\
$f_c$ & תדר מרכזי & 40 קילו-הרץ \\
$B$ & רוחב פס & 10 קילו-הרץ \\
$\lambda$ & אורך גל & $c/f_c$ מטרים \\
$r$ & טווח & מטרים \\
$v_r$ & מהירות רדיאלית & מטרים/שנייה \\
$\Delta t$ & זמן מעבר & שניות \\
$\text{SNR}$ & יחס אות לרעש & דציבלים \\
$P_{\text{total}}$ & צריכת הספק כוללת & 95 מיליוואט \\
$N$ & מספר חלקיקים או ריצות & 50 \\
$\omega$ & משקל מיזוג \en{Covariance Intersection} & $[0,1]$ \\
$\tau_m$ & קבוע זמן ממברני & 10 מילישניות \\
$V_{\text{th}}$ & סף פוטנציאל ממברנה & מיליוולטים \\
\bottomrule
\end{tabular}%
}
\end{table}